% Page 053

\seite{53}

das Schuljahr 1915/16\ob
Entlassen wurde kein Schüler, aufgenommen wurden 5 Mäd\ohb
chen. Die Schülerzahl beträgt 32 – 13 Knaben, 19 Mädchen. Alle\ob
Kinder sind katholisch.

\pend
\abschnitt{Ernte}
\pstart

Die Gemeinde erntete in diesem Jahre eine mittlere Ernte. Weizen\ob
Gerste und Hafer gediehen gut. Die Kartoffelernte ist zufriedenstellend.\ob
Die Wiesenpreise waren den Sommer über sehr hoch, besonders die\ob
Butter, die Preise der Lebensmittel waren gegen den vorstehenden\ob
Jahre sehr viel höher. Butter und Eier kosteten 1,50 – 1,60 M, während Eier die\ob
Preise 1,00 – 1,20 M betrugen.

\pend
\abschnitt{Fortsetzung}
\pstart

Während unsere Soldaten im Osten und Westen, zu Land und\ob
zur See mit bewunderungswürdiger Tapferkeit die Truppen des\ob
des Feindes ertragen, und so das Land von Feind und Hunger\ob
forthalten, wird auch die Bevölkerung aufs neue, weitere\ob
schwere Opfer ihres Vaterlandes zu bringen. Zunächst steht auf,\ob
die Obligation geht zurück. Die besten Rohstoffe sind\ob
längst eingelagen. Landes- und Kreisausschüsse, die eingesetzt haben,\ob
hat die Arbeit vielfach geleistet, fort. Kinder und alte Leute\ob
helfen sich fleißige Arbeit dem lieben Vaterlande nütz\ohb
lich zu erweisen, und in diesem Sinne werden die Kinder\ob
der Christlichen Bürgerschule fort vom Schuljahre befreit, außerdem\ob
die übrigen Klassen, daß Anordnung der Behörden ein Halb\ohb
tagsschule im Sommer halten. Zu wiederholten Male werden\ob
in den Volksschulen frei von Krieg für den Vaterländischen\ob
Jahresbericht die Schulkinder erfahren jede Vergünstigungen\ob
für Kriegswaisen. Es sind fünf im Dorfe 4 Krieger gestellt,\ob
ein Soldat ist seit einem Jahre vermißt. Die ganze Pfarrei
\pend
